\section{Method}

\subsection{Tokenizer Extension}

We fix the tokenizer to Glot500-style \emph{append-only injection}:
\begin{enumerate}
  \item Copy XLM-R's \texttt{sentencepiece.bpe.model} (250,002 tokens) as the source.
  \item Train a SentencePiece \emph{unigram} auxiliary model on the 99-language mixed
        corpus ($\alpha=0.3$ sampling) with \texttt{byte\_fallback=true} and
        \texttt{character\_coverage=1.0}.
  \item Select pieces \emph{absent} from the source SPM and append them directly to the
        SPM protobuf, preserving unigram scores.
        We do not use HuggingFace \texttt{add\_tokens()} to stay closer to the Glot500
        reference implementation.
\end{enumerate}

The resulting vocabulary expands from 250,002 to 366,666 tokens (116,664 new pieces).
Existing token IDs are preserved, but the \verb|<mask>| special token shifts from ID 250001
to ID 366665; we remap it by string identity during initialization (\S\ref{subsec:init}).

\paragraph{Difference from Yamaguchi et al.}
Yamaguchi et al.\ \citep{yamaguchi2025vocab} train the auxiliary tokenizer on a small
target-only corpus ($\approx$0.01~GB) and append only target-specific pieces, keeping
the vocabulary increment small.
Our Glot500-style approach trains on 99 mixed language-scripts and appends all 116,664
new pieces.
These two designs differ in both training data scope and vocabulary increment size;
we fix this axis to Glot500-style and vary only initialization.

\subsection{Embedding Initialization}
\label{subsec:init}

Three rules apply to all five methods:
(i) source-tokenizer tokens copy their source row exactly;
(ii) only the 116,664 new rows are method-specific;
(iii) the \verb|<mask>| row is remapped by string identity, and input embedding and
LM head weight tying is maintained.
Audit numbers: source length 250,002; target length 366,666; new rows 116,664;
\verb|<mask>| copy diff 0.0; LM head tied.

\begin{table}[h]\centering\small
\caption{Five initialization methods for new embedding rows.}
\label{tab:methods}
\begin{tabular}{@{}lp{9cm}@{}}
\toprule
Method & Initialization of new row $\mathbf{e}_j$ \\
\midrule
\rand & $\mathbf{e}_j \sim \mathcal{N}(0,\,0.02^2\mathbf{I})$ \\
\meanm & $\mathbf{e}_j = \mathrm{mean}(\mathbf{E}_\mathrm{src})$ \\
\fvt & $S = \mathrm{SrcTok}(\mathrm{surf}(t_j))$; $\mathbf{e}_j = \mathrm{mean}(\mathbf{E}_\mathrm{src}[S])$ \\
\wfvt & $w_s = |\mathrm{piece}(s)|$; $\mathbf{e}_j = \big(\sum_s w_s \mathbf{E}_\mathrm{src}[s]\big)/\sum_s w_s$ \\
\famm & $f = \mathrm{family}(\mathrm{prov}(t_j))$; $\mathbf{e}_j = \mathrm{mean}(\mathbf{E}_\mathrm{src}[\mathrm{SrcTok}(f)])$ \\
\bottomrule
\end{tabular}
\end{table}

\textbf{\rand} draws from $\mathcal{N}(0,0.02)$, providing no information.
\textbf{\meanm} starts at the global centroid of pretrained embeddings---a stable but
uninformative starting point.
\textbf{\fvt} decomposes each new token's surface form with the source tokenizer and
averages the constituent source embeddings, reusing surface/subword information
(e.g., a new token \texttt{\_bambara} decomposed by XLM-R into
\texttt{[\_bam][bara]} starts at the mean of those two rows).
\textbf{\wfvt} weights this average by surface-piece length, giving more weight to
longer (more specific) source subword pieces.
\textbf{\famm} replaces surface-form decomposition with a typological prior: it aggregates
provenance information from the tokenizer training corpus to determine which language-script
each new token most often appears in, maps that language-script to its family, and
initializes the new row at the centroid of source embeddings belonging to head languages
in the same family.

\paragraph{Family assignment for \famm.}
Table~\ref{tab:fam} shows which families Target7 languages belong to and which head
source languages \famm\ can use.
For \texttt{dtp} (Austronesian), the prior draws from \texttt{ind/msa/jav/tgl/mlg}.
For \texttt{bam} (Mande) and \texttt{xav} (Macro-Je), no head language belongs to the
same family; \texttt{bam} falls back to the nearest Niger-Congo head (\texttt{xho}),
and \texttt{xav} has no usable family prior.

\begin{table}[h]\centering\small
\caption{Target7 family assignments and head source languages used by \famm.}
\label{tab:fam}
\begin{tabular}{@{}lll@{}}
\toprule
Target & Family & Head source languages \\
\midrule
\texttt{dtp} & Austronesian & \texttt{ind, msa, jav, tgl, mlg} \\
\texttt{csb} & Slavic & \texttt{ces, hrv, pol, slk, slv} \\
\texttt{fur} & Romance & \texttt{cat, fra, ita, por, ron} \\
\texttt{lij} & Romance & \texttt{cat, fra, ita, por, ron} \\
\texttt{ile} & Constructed & \texttt{epo} \\
\texttt{bam} & Mande & Niger-Congo proxy: \texttt{xho} \\
\texttt{xav} & Macro-Je & none \\
\bottomrule
\end{tabular}
\end{table}

\subsection{Continued MLM Pretraining}

We continue pretraining the initialized model with the MLM objective on the corpus
described in \S3.
All five methods use identical hyperparameters; only the starting checkpoint differs.

\begin{table}[h]\centering\small
\caption{Continued MLM pretraining settings.}
\label{tab:mlm}
\begin{tabular}{@{}ll@{}}
\toprule
Item & Value \\
\midrule
Base model & \texttt{xlm-roberta-base} \\
Objective & MLM (dynamic masking, mask prob. 0.15) \\
Learning rate & $5\times10^{-5}$ \\
Adam & $\beta=(0.9,0.999)$, $\varepsilon=10^{-8}$ \\
LR schedule / warmup / weight decay & linear / 0 / 0.0 \\
Max sequence length & 512 \\
Language sampling & $\alpha=0.3$ \\
Checkpoint interval & 1,000 steps \\
Convergence budget & 50,000 steps (all five methods) \\
\bottomrule
\end{tabular}
\end{table}
